% =====================================================================
%  Conclusiones
% =====================================================================
\chapter[CONCLUSIONES]{Conclusiones}
\label{cap:conclusiones}

Durante el Trabajo Terminal I se estableció el marco referencial del proyecto donde se revisaron los
fundamentos del EEG, la componente P300 y el paradigma P300 Speller, además de las técnicas de
preprocesamiento para usar estas señales, junto a las arquitecturas de aprendizaje profundo y las
métricas con las que se evalúan estos sistemas. Por otro lado el estado del arte mostró que la
mayoría de los sistemas P300 Speller basados en EEG alcanzan una alta exactitud en entornos
controlados, pero bajo tres condiciones que dificultan su uso fuera del laboratorio. La primera es
el registro con muchos canales, que en la base de datos de la competencia BCI III llega a 64. La
segunda es el alto costo de los equipos necesarios para registrarlos.

A partir de lo anterior se desarrolló una propuesta con dos partes. La primera es un modelo de
aprendizaje profundo para detectar la componente P300 y la segunda es el sistema de comunicación
asistida que integra ese modelo junto con un módulo de predicción de palabras para agilizar la
construcción de palabras. Para cada parte se planteó una metodología, CRISP-ML para el modelo e
incremental para el sistema. Ambas empatan en el último incremento, cuando la arquitectura
seleccionada se incorpora a la aplicación para entrenar el modelo de cada usuario al terminar su
calibración.

En el Trabajo Terminal II quedan por construir los tres incrementos del sistema. El primero es el
módulo de predicción de palabras, cuya técnica aún está por definirse. Le siguen la interfaz con la
matriz del P300 Speller y, al final, la calibración junto con la integración del modelo de
detección. En paralelo se completarán las fases de CRISP-ML, con el entrenamiento de las
arquitecturas candidatas, y el ajuste de sus hiperparámetros. También se registrarán las señales
propias de voluntarios, con las que se probará el sistema completo entre abril y mayo.

Queda pendiente verificar si el hardware seleccionado es compatible con un microordenador NVIDIA
Jetson, en el que se analizaría la ejecución de la inferencia del modelo. Si el casco no puede
comunicarse con esta tarjeta, la Jetson no se usará y el sistema se implementará únicamente en la
computadora con Windows, sin que esto afecte el resto del alcance.
